\chapter{Conclusiones generales}

\section{Aportes a la caracterización de impactos climáticos}

Esta tesis aporta una ruta integrada para caracterizar impactos de variabilidad climática global sobre el rendimiento de cultivos básicos. Al considerar ENSO y MJO de manera conjunta, el trabajo permite conectar escalas interanuales e intraestacionales del clima con anomalías de rendimiento agrícola. Esta aproximación amplía la lectura tradicional centrada exclusivamente en ENSO y reconoce que eventos de menor escala temporal pueden coincidir con ventanas fenológicas críticas.

% TODO: completar con los hallazgos empíricos finales por cultivo y región.

\section{Aportes a la identificación de regímenes agroclimáticos}

La tesis propone la construcción de regímenes agroclimáticos interpretables que integran sensibilidad climática, ubicación geográfica, cultivo y fase fenológica. Estos regímenes pueden servir como una síntesis útil para comparar regiones agrícolas globales, identificar patrones de vulnerabilidad y orientar análisis posteriores de riesgo.

% TODO: completar con número, nombre y descripción final de los regímenes identificados.

\section{Aportes a la generación de insumos para seguros paramétricos}

El trabajo traduce relaciones clima-rendimiento en insumos preliminares para diseño de seguros paramétricos agrícolas. En particular, propone identificar umbrales climáticos, ventanas temporales y regiones prioritarias donde la relación entre índices climáticos y anomalías negativas de rendimiento pueda ser suficientemente clara para explorar disparadores paramétricos.

% TODO: completar con triggers finales, desempeño preliminar y regiones de mayor potencial.

\section{Limitaciones metodológicas}

Las principales limitaciones metodológicas se relacionan con la calidad y resolución de los datos de rendimiento, la incertidumbre de los calendarios fenológicos globales, la heterogeneidad del manejo agrícola, la presencia de riego, la variabilidad de variedades cultivadas y la dificultad de separar efectos climáticos de tendencias tecnológicas o socioeconómicas. Además, las relaciones estadísticas entre índices climáticos y rendimiento no deben interpretarse automáticamente como causalidad directa.

% TODO: completar con limitaciones específicas encontradas durante el análisis.

\section{Potencial para una herramienta web o sistema de consulta}

Los productos de la tesis pueden llevarse a una herramienta web o sistema de consulta que permita explorar mapas de sensibilidad, regímenes agroclimáticos, ventanas fenológicas críticas y triggers paramétricos preliminares por cultivo y región. Una herramienta de este tipo podría apoyar análisis de riesgo, priorización territorial, diseño de productos financieros y comunicación entre investigadores, aseguradoras, gobiernos y entidades de seguridad alimentaria.

% TODO: definir arquitectura preliminar, capas de información y usuarios objetivo de la herramienta.
